% s-068-p068-m001; original PDF 68, 69, 70
\Aphorism{s-068-p068-m001}{63}
Wenn ich bedenke, wie erfinderisch und tüchtig der Mensch im Ausbau seiner eigenen Lebensführung sein kann, mit welcher neapolitanischen Anmuth er gegenüber Zufall und Mitmenschen sich durchzusetzen weiss, ohne Grossthuerei und moralische Draperie, ohne den schielenden Seitenblick auf Andere und die „Wohlfahrt Anderer“ — und wie plump hingegen, wie anmassend und effectlüstern sich der Altruist in fremde Schicksale eindrängt und dazwischenschiebt, immer auf der Jagd nach moralischer Selbstbefriedigung und dem Kitzel des „guten Gewissens“; wenn ich mir diese beiden Typen gegenüberstelle, den Egoisten \Foreign{optima fide}, der aus seiner inneren Unschuld noch an Andere abgiebt, indem er ihnen zu einem pharisäischen Distanzbewusstsein von „jenem Zöllner“ verhilft, und den selbstlosen Unschuldsaspiranten, der auf Hintertreppen zu seiner Höhe steigt und nur desshalb steigt, um die Untenbleibenden zu beschämen: dann zögere ich nicht einen Augenblick, als Mensch von \Sperr{Geschmack} für den Egoisten zu sein, selbst wenn ich als Mensch von Moral für den Andern sein müsste. Auch die Selbstliebe hat ja ihre Grade von Vornehmheit, und vielleicht stehen ihre untersten Grade tiefer als der Altruismus selbst in seinen gemeinsten Äusserungsformen jemals sinken kann, sodass es in Hinsicht auf die grosse Masse, die von beiden Gattungen nur die niedrigsten Arten darzustellen vermag, in der That geboten sein wird, den moralischen Menschen vor dem egoistischen \Foreign{a priori} zu bevorzugen. Dafür aber stehen die eigentlich hohen Typen des Individualismus unvergleichlich viel höher als der moralische Mensch auch in seiner edelsten Gestalt, und ich für meine Person sehe in der \Sperr{Ehrfurcht vor sich selbst als einem Schaffenden} das höchste Ethos, das auf Erden erreichbar ist: gegen diesen glaubenssicheren Gnadenstand ist die kleinlich aufgeregte Werkheiligkeit der moralischen Menschen beinahe etwas Verächtliches. Der moralische Mensch schleicht sich durch eine Hinterthür in den Seelenzustand des naiv-schöpferischen Menschen ein; dieser ist gerechtfertigt, und jener liegt auf der Lauer nach Rechtfertigung, dieser hat Unschuld als Grundgeheimniss seines Wesens, jener angelt nach ihr mit Äusserlichkeiten und Gelegenheitsursachen; dieser setzt sie bei sich voraus, jener muss sie sich von aussen her, durch Handlungen und Handgreiflichkeiten, beweisen! — Aber gerathen wir nicht selbst ins Moralisiren hinein; bleiben wir beim \Foreign{sic volo, sic jubeo} des Geschmacks. Es ist altruistischer Geschmack, dem „nur das Gute wohlthut, das wir bei offenen Thüren geniessen können und das allen Menschen nützt“ (Emerson); es ist unser Geschmack, der sich die Frage stellt: kann eine Species von Menschen und Handlungen überhaupt etwas für sich Werthvolles sein, die nur bei offenen Thüren gedeiht, ja die nur \Foreign{sub specie} der Öffentlichkeit \Sperr{existirt}? Der moralische Mensch und die moralische Handlung sind ja nur \Sperr{relative} Realitäten, da sie als nothwendige \Foreign{ratio essendi} ein Forum moralischer Beurtheilung voraussetzen; sie beruhen, Kantisch gesprochen, auf Heteronomie, nicht auf Autonomie! So schön man auch von der Selbstgenugsamkeit der Tugend und der immanenten Moralgerichtsbarkeit des Gewissens gefabelt hat: der Historiker der moralischen Empfindungen weiss zu genau, dass diese immanente Sittlichkeit im Grunde nur eine gesellschaftlich geforderte Sitte ist, und dass der moralische Mensch nie sich selbst beurtheilt, sondern stets von wirklichen oder vorgestellten „Anderen“ beurtheilen lässt. Ist eben dadurch der moralische Mensch nicht a priori von der höchsten Vornehmheit ausgeschlossen? Wirkt die altruistische Lebensweise, dieses grundsätzliche „Leben für Andere“ und Rücksichtnehmen auf Andere, das beständige Offenhalten aller Thüren und Sich-in-Scene-setzen vor Zuschauern und Kritikern — wirkt dieses ewige Coulissenleben nicht jeder ruhigen und einsiedlerischen Seelenbildung entgegen? Muss ein Mensch, der nur das Licht fremder Meinungen und Werthschätzungen reflectirt und sofort zum eigenschaftslosen Massenklumpen verlöschen würde, sobald es der Aussenwelt einmal einfiele, jene Werthschätzungen einzustellen, — muss solch ein hypothetischer Mensch, der sich nur ex concessis beweisen kann, gegen den schöpferisch-kategorischen Menschen nicht eine erbärmliche Figur abgeben? Übrigens, und darin liegt der Humor, ist auch der moralische Mensch \Sperr{Egoist}, nur dass er seinen Ernst und Idealismus nicht dem ächten, intim bekannten, vor Daseinsrausch productiven Vollblut-Ich zuwendet, sondern einem an der Umwelt reflectirten Denk- und Begriffs-Ich, das den Hin- und Rückweg durch ein absorbirendes Medium genommen und sich von jenem ursprünglichen Ich nicht um eine, sondern gleich um zwei Helligkeitsstufen entfernt hat.
\par

